Determina si las siguientes tablas de datos son o no una relación proporcional:

        \begin{multicols}{2}
            \begin{parts}
                \part
                \begin{minipage}[t]{0.35\linewidth}
                    \includegraphics[width=\textwidth]{mex_0074.png}
                \end{minipage}%
                \begin{minipage}[b]{0.6\linewidth}
                    \begin{solutionbox}{3.5cm}\small%

                        \begin{multicols}{2}
                            $1\div  6=\frac{1}{6}$\\[0.3em]
                            $2\div 12=\frac{1}{6}$\\[0.3em]
                            $3\div 18=\frac{1}{6}$\\[0.3em]
                            $4\div 24=\frac{1}{6}$\\[0.3em]
                            $5\div 30=\frac{1}{6}$

                            \columnbreak% 

                            $\therefore$ La constante de proporcionalidad es $\dfrac{1}{6}$.                        
                        \end{multicols}
                    \end{solutionbox}
                \end{minipage}

                \part
                \begin{minipage}[t]{0.35\linewidth}
                    \includegraphics[width=\textwidth]{mex_0075.png}
                \end{minipage}%
                \begin{minipage}[b]{0.6\linewidth}
                    \begin{solutionbox}{3.5cm}\small%

                        \begin{multicols}{2}
                            $\frac{3}{4}\div  1=\frac{3}{4}$\\[0.3em]
                            $\frac{3}{2}\div  2=\frac{3}{4}$\\[0.3em]
                            $\frac{9}{4}\div 3=\frac{3}{4}$\\[0.3em]
                            $3\div 4=\frac{3}{4}$\\[0.3em]
                            $\frac{15}{4}\div 5=\frac{3}{4}$

                            \columnbreak%  

                            $\therefore$ La constante de proporcionalidad es $\dfrac{3}{4}$.
                        \end{multicols}
                    \end{solutionbox}
                \end{minipage}

                \part
                \begin{minipage}[t]{0.35\linewidth}
                    \includegraphics[width=\textwidth]{mex_0076.png}
                \end{minipage}%
                \begin{minipage}[b]{0.6\linewidth}
                    \begin{solutionbox}{3.5cm}\small%

                        \begin{multicols}{2}
                            $\frac{6}{5}\div  1=\frac{6}{5}$\\[0.3em]
                            $\frac{12}{5}\div 2=\frac{6}{5}$\\[0.3em]
                            $\frac{18}{5}\div 3=\frac{6}{5}$\\[0.3em]
                            $\frac{24}{5}\div 4=\frac{6}{5}$\\[0.3em]
                            $6\div 5=\frac{6}{5}$

                            \columnbreak%  

                            $\therefore$ La constante de proporcionalidad es $\dfrac{6}{5}$.
                        \end{multicols}
                    \end{solutionbox}
                \end{minipage}

                \part
                \begin{minipage}[t]{0.35\linewidth}
                    \includegraphics[width=\textwidth]{mex_0077.png}
                \end{minipage}%
                \begin{minipage}[b]{0.6\linewidth}
                    \begin{solutionbox}{2.8cm}\small%

                        \begin{multicols}{2}
                            $ 8 \div 1=8$\\[-0.2em]
                            $ 16\div 2=8$\\[-0.2em]
                            $ 24\div 3=8$\\[-0.2em]
                            $ 32\div 4=8$\\[-0.2em]
                            $ 40\div 5=8$

                            \columnbreak% 

                            $\therefore$ La constante de proporcionalidad es 8.
                        \end{multicols}
                    \end{solutionbox}
                \end{minipage}
            \end{parts}
        \end{multicols}
    